\documentclass[11pt]{article}
\usepackage[T1]{fontenc}
\usepackage{lmodern,microtype,amsmath,amssymb,amsthm,booktabs,tabularx,array}
\usepackage[letterpaper,margin=1in,top=.9in,bottom=.92in]{geometry}
\usepackage{xcolor,fancyhdr,enumitem}
\definecolor{inkblue}{HTML}{193E56}
\usepackage[colorlinks=true,allcolors=inkblue,bookmarksnumbered=true,pdfauthor={Research analysis},pdftitle={Prime-exponent patterns and p-adic branch precision}]{hyperref}
\newtheorem{theorem}{Theorem}[section]
\newcommand{\upstream}[2]{\href{https://github.com/lsernau/grh-lean/blob/72fead102f663d5c7a407964a13dc0cf9d977207/#1}{#2}}
\newcommand{\incoming}[2]{\href{https://github.com/VladimirReshetnikov/ProveIt/blob/1677488425e45c3df14149c804014133798136a5/docs/incoming/#1}{#2}}
\pagestyle{fancy}\fancyhf{}\lhead{\small\textcolor{inkblue}{Prime-exponent patterns and branch precision}}\rhead{\small Research note}\cfoot{\thepage}
\setlength{\headheight}{15pt}\setlength{\emergencystretch}{2.5em}
\setlist{nosep}\setcounter{tocdepth}{1}
\begin{document}
\thispagestyle{plain}
{\small\sffamily\color{inkblue} MATHEMATICAL RESEARCH NOTE\hfill 26 NOVEMBER 2026}\par\vspace{.5em}
{\Huge\bfseries Prime-exponent patterns and\par $p$-adic branch precision}\par\vspace{.6em}
{\Large Two transfers from the Dirichlet GRH development}\par\vspace{.7em}
\noindent\rule{\linewidth}{.65pt}
\paragraph{Results.}
McNew conjectured that Rankin's set has the greatest natural density among sets of integers avoiding three-term geometric progressions of rational ratio. An incoming ProveIt manuscript supplied a counterexample. We give a different proof that changes only the lowest ternary digit of each selected prime exponent. Removing two vertices of a binary cube suffices; we classify every profitable replacement with at most two removals. In a degree-ten number field a twelve-prime example gives an exact relative gain of $1/11337408$. For the integers, the supremal density under arbitrary independent rules on blocks of at most $d$ primes is attained, computable, and strictly below the unrestricted optimum for every fixed $d$. These optima tend to the unrestricted optimum. For every block size up through four, the optimum equals Rankin's density exactly.

A second question in the incoming collection asks when a pole bound is sharp in reconstructing a Taylor polynomial from finitely many Lubin--Tate branches. We give a necessary and sufficient base-$p$ carry test, prove eventual periodicity of the bounded error about its linear growth, and determine the optimal sampling directions at every modulus. Over $\mathbb Q_p$ the result gives the exact simultaneous vanishing threshold for a family of weighted binomial sums and the full vector of residues one step before vanishing.

\paragraph{Scope and provenance.}
The two arguments use internal methods of the Lean project and earlier mathematical constructions. Their proofs require no hypothesis about zeros. We read pinned source and selected incoming manuscripts; the accompanying Python checks reproduce finite cases independently of Lean. The upstream authors report kernel checks of their five exported theorems. This study did not rebuild their project or formalize the extensions. We identify antecedents below; historical priority has not been established for the refinements.

\tableofcontents
\newpage
\section{From the formal proof to two concrete questions}
The readable edition of \texttt{lsernau/grh-lean} at commit \texttt{72fead102f66} ends in five theorems. They give mathlib's Riemann hypothesis, an explicit zeta zero-line statement, nonvanishing of every Dirichlet $L$-function for $\Re s>1/2$ with the principal pole excepted, a zero-line statement for its open critical strip, and the primitive-character specialization. The modulus may be any positive integer, and the statement includes imprimitive characters. The statements use Dirichlet characters over $\mathbb Q$. The proof is attributed upstream to Levent Alp\"oge's work with Claude, with a shorter formal development by Sernau and collaborators. See the \upstream{README.md}{repository description}, \upstream{grh/src/Statements.lean}{definitions}, \upstream{grh/src/Main.lean}{five exports}, and \upstream{grh/README.md}{verification record}.

The mechanism behind the zero theorem matters for deciding what transfers. The proof constructs integral automorphic forms and estimates a lattice from two directions. Short forms supply upper bounds, a second moment ensures sufficiently many independent images, and a lower bound for the determinant restricts how short the forms can simultaneously be. In the readable edition the comparison yields
\[
        m\,\mathrm{totalRemainder}\ \le\ f\,a\,(\log b)^5.
\]
Here $m$ is the number of lifts, $a$ is a derivative order and $f$ counts infinite places. The proof takes $m$ of order $\sqrt{x}$ and can take $\log b\le2x^\epsilon$. After normalization and control of the auxiliary scale it bounds a smoothed, centered prime sum from above at order $x^{-1/2+6\epsilon}$ for integer $x$ with fixed arithmetic data. The sum uses $\Lambda(n)(1-\chi_K(n))$ in a residue class. A one-sided analytic criterion converts these bounds to zero exclusion. See the \upstream{grh/GUIDE.md}{reading guide}, \upstream{grh/src/New/Reduction/MainInequality.lean}{main inequality}, and \upstream{grh/src/New/Reduction/OneSidedBound.lean}{one-sided estimate}.

Two smaller components admit independent tests.
\begin{itemize}
\item In \upstream{grh/src/New/Reduction/Data/Existence.lean}{\texttt{exists\_three\_term\_free}}, base expansion and a sphere give a large set avoiding weighted linear relations. A later lemma uses exponents to turn simultaneous multiplicative restrictions on places into additive restrictions. The source credits Behrend and Shapira. Our density argument also passes from multiplicative progressions to digit vectors. Its new ingredient is a weighted modification of a cap and an exact analysis of its secants; it uses neither the sphere count nor the tower parameters.
\item In \upstream{grh/src/New/Derivatives/Sharpness/OneCurve.lean}{\texttt{OneCurve.lean}}, ultrametric bounds on coefficients are attained by a selected monomial. The one-curve derivative extremum is attained at a power of a local coordinate. The incoming interpolation report applied the idea of exhibiting an extremizer to intersecting branches. We compute the extremal pole after cancellation modulo every $\pi^r$. The extremizer here is a polynomial vanishing on the other branches.
\end{itemize}
\paragraph{A limit of direct analytic transfer.}
Take the splitting field of $x^3-x-1$. The polynomial is irreducible by the rational-root test and has discriminant $-23$, so its Galois group is $S_3$. At unramified primes the identity and a three-cycle have the same sign, while their traces in the two-dimensional standard representation are $2$ and $-1$. In fact, for any cyclotomic field $C$, its intersection with the splitting field is abelian and lies in the quadratic subfield. Both even classes can therefore lift with identity on $C$ in the compositum. Chebotarev gives primes of each type with this identical congruence information. A proof for the degree-two Artin factor would need arithmetic estimates that see this difference and a compatible determinant comparison. The terminal zero criterion alone supplies neither estimate. The cubic example was already highlighted in the incoming reports \cite{priorideal}.
\section{A counterexample using one ternary digit}\label{sec:cap}
A set of positive integers avoids rational-ratio geometric progressions if $ac=b^2$ with $a,b,c$ in the set forces $a=b=c$. Let $J$ be the nonnegative integers whose base-three digits are all zero or one. Rankin's set $R$ consists of the integers with $v_p(n)\in J$ for every prime $p$. The equation $a+c=2b$ in $J$ forces equality: neither side has carries, and each digit of $2b$ is zero or two. Applying this at each prime proves the required property of $R$. McNew's Conjecture 7.2 asserts that its natural density, approximately $0.71974$, is optimal \cite{mcnew}. A manuscript already in the incoming collection disproves this by altering exact valuations at finitely many primes \cite{incomingmc}. Here the modification is periodic in the lowest ternary digits. This permits an exact classification of the smallest possible number of deleted digit patterns.

\subsection{The construction}
Fix a finite list of distinct primes, one called the marker and $k$ called satellites. Write $w$ for the reciprocal of the marker and $u_1,\ldots,u_k$ for the reciprocal satellite primes. A \emph{cap} is a subset of $\mathbb F_3^{k+1}$ in which $x+y+z=0$ implies $x=y=z$. Use representatives $0,1,2$. The binary cube $C=\{0,1\}^{k+1}$ is a cap. Write $b=(1,\ldots,1)$ for the satellite vector, and $b_i$ for the vector obtained by replacing its $i$th entry by 2. Form
\[
 C'=\bigl(C\setminus\{(0,b),(1,b)\}\bigr)\ \cup \{(2,b),(2,b_1),\ldots,(2,b_k)\}.
\]
\textbf{Claim.} $C'$ is a cap. Three old points lie in $C$. Two new points force the marker of the third to equal 2; all three then have satellite digits in $\{1,2\}$, which is itself a cap in each coordinate. If there is exactly one new point, every satellite coordinate where it equals 1 forces both old digits to be 1. At its possible exceptional position, the old digits must be 0 and 1. Thus at least one old point has satellite vector $b$, and both such old points were removed. This covers all configurations.

Define $E$ as follows. At the selected primes write the exponent vector as $3h+c$, with each coordinate of $c$ in $\{0,1,2\}$. Require $c\in C'$ and $h\in J^{k+1}$. Use Rankin's rule at all other primes. In a progression the three residue vectors sum to zero modulo three, even when some valuation steps are negative. The cap property makes them equal. Subtract the common vector, divide by three and use the progression-freeness of $J$ coordinate by coordinate. The unselected valuations are constant for the same reason. Hence $E$ avoids every nonconstant rational-ratio progression.

Give a digit vector $c$ mass $W(c)=w^{c_0}\prod_i u_i^{c_i}$. The exact unnormalized mass change is
\begin{equation}\label{eq:gain}
 W(C')-W(C)=\left(\prod_i u_i\right)
 \left[w^2\left(1+\sum_i u_i\right)-(1+w)\right].
\end{equation}
Here $W$ on sets denotes a sum. To obtain the change in natural density multiply by the positive factor
\[
 \prod_{p\in S}\left[(1-p^{-1})\sum_{h\in J}p^{-3h}\right]
 \prod_{p\notin S}\left[(1-p^{-1})\sum_{e\in J}p^{-e}\right].
\]
Indeed, at finitely many primes the Chinese remainder theorem gives the independent geometric laws $(1-1/p)p^{-e}$. Large exponents at any fixed prime have a vanishing geometric tail. Outside $S$, a violation at $p>Y$ requires $p^2\mid n$, costing at most $\sum_{p>Y}p^{-2}$ uniformly in the counting cutoff. The infinite product is positive because each full local factor is at least $1-p^{-2}$.

\begin{theorem}\label{th:rankin}
Use marker $2$ and take the first odd primes until their reciprocals sum to more than $5$. Then $d(E)>d(R)$. Only the units digits in base three at finitely many primes differ from Rankin's rule.
\end{theorem}
Euler's divergence of $\sum_p1/p$ ensures termination. With $w=1/2$, the bracket in \eqref{eq:gain} is $(\sum_i u_i-5)/4$. This proves strict gain in the ordinary limiting density and refutes the published conjecture. One may even take all odd primes up to $Y=\lceil\exp(\exp(7))\rceil$. The finite Euler product over $p\le Y$ exceeds $\sum_{n\le Y}1/n\ge\log Y$. In its logarithm the contribution of prime powers of degree at least two is bounded by $\sum_{n\ge2}1/(n(n-1))=1$. Subtracting the contribution $1/2$ of the prime 2 leaves at least $7-1-1/2>5$. The earlier certificate in \cite{incomingmc} uses only 224 odd primes up to 1427 and an exact-valuation rule. The point of the digit construction is its short proof and its optimality under the following restriction.

\section{Exactly which two deletions can help?}\label{sec:classification}
Let $C=\{0,1\}^n\subset\mathbb F_3^n$ carry the weights $W(x)=\prod_i w_i^{x_i}$, with $0<w_i\le1/2$. Require that a replacement cap keep all cube vertices except a specified set $D$ and omit every vertex of $D$. The question is how much outside weight it can add.

\begin{theorem}\label{th:edge}
Deleting zero or one vertex, or two vertices at Hamming distance at least two, cannot yield a net gain. For two adjacent deleted vertices, call their varying coordinate the marker, of weight $w$, and their common remaining vector $b$. For a satellite $i$ put $s_i=w_i^{2-b_i}$. The maximum possible added weight equals
\[
 W(b)\max\!\left\{(1+w)\max_i s_i,\;
               w^2\Big(1+\sum_i s_i\Big)\right\},\qquad
 W(b)=\prod_{i\ne\mathrm{marker}} w_i^{b_i}.
\]
Use zero for an empty maximum. The deleted weight equals $(1+w)W(b)$.
\end{theorem}

\begin{proof}
An outside point with $h$ entries equal to 2 lies on exactly $2^{h-1}$ lines through pairs of cube points. At a 2 the two binary digits are opposite, and everywhere else both equal the digit of the outside point. These pairs are disjoint. Thus $h\le2$ after two deletions. For $h=2$, both deletions must belong to its square face and must meet different opposite pairs. In a binary square this means that they are adjacent. For $h=1$ at least one endpoint of the associated cube edge must have been deleted.

For a single missing vertex $v$, the allowed additions therefore arise by changing a single digit of $v$ to 2. Two different such changes have as third point the cube vertex obtained by flipping both of those binary digits in $v$. It survives. Thus at most one point can be added, of weight at most $W(v)/2$. For two nonadjacent missing vertices, each possible addition is incident to exactly one of them. Two different additions incident to $v$ require that their two-digit flip equal the other missing vertex. At most two additions can therefore be incident to $v$. Each weighs at most $W(v)/2$, which proves the claim by summing charges. With no deletion the secant argument excludes every outside point.

Now delete an edge. The available points consist of the center $(2,b)$ and, for each $i$, the three marker choices above the satellite vector $b_i$ obtained by changing its $i$th digit to 2. Call that triple column $i$. The preceding face argument gives this entire list and shows that each individual candidate is compatible with all pairs of retained points.

Suppose different columns $i,j$ use marker digits $a,c$. Their satellite completion has $b$ flipped at $i$ and $j$ and is binary. Its marker $-a-c$ must be 2 or else it is a retained cube vertex. Thus $a+c=1$ modulo three. A center and a column point give a one-flip binary completion by the same calculation, forcing the column marker to be 2. Three or more occupied columns force every marker to be 2. Two occupied columns without the center can have only the pairs $(0,1),(1,0),(2,2)$, each column containing a single point. In one occupied column there are at most two points, of combined relative marker weight at most $1+w$. The two mixed-marker cases also cost at most $(1+w)\max_i s_i$ in units of $W(b)$. All cases with only marker 2 cost at most the second term of the maximum.

Both bounds are attained whenever the corresponding candidates exist. In a maximizing column use markers 0 and 1. Their completion is the third point in that same column and is absent. Alternatively use the center and all columns with marker 2. The satellite digits then belong to the cap $\prod_i\{b_i,2\}$. With only one new point, a pair of retained points cannot complete it. This proves the formula, including the no-satellite case.
\end{proof}
Since $s_i\le w_i$, for a fixed marker the most favorable sign of the change occurs at the all-ones satellite vector. The marker giving the easiest threshold is one of largest weight. Indeed with $S_0=\sum_iw_i$ fixed, the ratio $w^2(1+S_0-w)/(1+w)$ increases for $0<w\le1/2$: its logarithmic derivative is $2/w-1/(1+S_0-w)-1/(1+w)>0$. Consequently the threshold of $5$ odd-prime reciprocal mass is sharp for a two-deletion digit exchange with marker 2. For a prime list containing 2, no choice of another marker can lower the threshold. Theorem~\ref{th:rankin} also shows that two is the smallest possible number of deleted cube vertices in a profitable one-digit exchange. The restriction concerns the residue cap, and permits arbitrarily many added patterns.

\subsection{Natural density for rational functions and fractional ideals}
The same proof works for monic polynomials over every $\mathbb F_Q$. A prime $P$ has weight $Q^{-\deg P}$. Among monics of degree exactly $n$, simultaneous divisibility by any fixed monic $D$ has probability $Q^{-\deg D}$ once $n\ge\deg D$. The finite valuation laws follow by inclusion--exclusion. An omitted high-degree square condition costs at most $\sum_{j>J_0}Q^jQ^{-2j}$. Thus the modified set has an exact-degree limit, which is also its degree-ball density. The reciprocal sum over irreducibles diverges: for $s>1$ the Euler product is $(1-Q^{1-s})^{-1}$, while the powers of exponent at least two contribute a bounded amount to its logarithm as $s\downarrow1$. Choose a linear marker and enough satellites to make $\sum_i u_i>Q^2+Q-1$. Equation~\eqref{eq:gain} gives a strict gain over the greedy set even for ratios in $\mathbb F_Q(x)\setminus\mathbb F_Q$. Allowing all nonzero scalar multiples preserves the result. Asada et al. asked for an example with greater upper density in Question~6.4 \cite{asada}; this construction has a limiting density. Their Conjecture~6.3 uses a different benchmark $m_Q$ for polynomial ratios.

For the nonzero integral ideals of any number field, put $w_{\mathfrak p}=1/N\mathfrak p$. A progression may have an arbitrary nonidentity fractional ideal as ratio as long as its three terms are integral. The equality of valuations covers this signed-step convention. Counting all ideals of norm $\le X$ gives $A_K(X)\sim c_KX$, $c_K>0$ \cite{idealcount}. Hence $A_K(X/NJ)/A_K(X)\to 1/NJ$ for any fixed ideal $J$, and the same quotient is $O_K(1/NJ)$ uniformly for $X$ large. This proves the geometric valuation laws and controls the square tail by $O_K(\sum_{N\mathfrak p>Y}(N\mathfrak p)^{-2})=o(1)$. The simple pole of the Dedekind zeta function at 1 and the bounded higher prime-power contribution again give $\sum_{\mathfrak p}(N\mathfrak p)^{-1}=\infty$. One finite edge exchange therefore beats the natural density of the greedy ideals in every field under this stronger ratio convention. The quadratic-field questions and the previous integral-ideal exchange are described in \cite{best,priorideal}.

\paragraph{A sharp small example for ideals.}
Within this two-deletion class, at least 12 prime coordinates are needed in any number field. In fact $w\le1/2$ and $\sum_i u_i> w^{-2}+w^{-1}-1\ge5$ requires at least 11 satellites. With exactly 12 coordinates the smallest possible field degree is 10. A field of degree at most 9 has at most nine norm-2 prime ideals, by the degree formula for primes over 2. If a norm-2 prime is present, use the most favorable marker of weight $1/2$. At most eight of the eleven satellites then have weight $1/2$, so their total is at most $8/2+3/3=5$. Without such a prime the bound is even smaller.

For equality of the dimension and degree thresholds, take the unique degree-10 subfield of $\mathbb Q(\zeta_{431})$. The integer 431 is prime, its cyclotomic Galois group is cyclic of order 430, and the kernel defining this subfield has order 43. Repeated squaring gives $2^{43}\equiv3^{43}\equiv1\pmod {431}$. Both primes therefore split completely. Take ten norm-2 primes and two norm-3 primes. For a norm-2 marker the eleven satellite weights sum to $9/2+2/3=31/6$. Formula~\eqref{eq:gain} gives an added unnormalized mass of $1/110592$. Dividing by the original cube mass $(3/2)^{10}(4/3)^2$ gives the exact relative result
\[
                 d(E)=d(R_K)\left(1+\frac1{11337408}\right).
\]
Here $R_K$ denotes the greedy ideal set of this field. The small degree assertion concerns the stated 12-coordinate, two-deletion problem.

There is also no field-independent upper bound for the ratio of the achievable density to the greedy ideal density. Choose $L$ distinct rational primes congruent to 1 modulo 8 and adjoin their square roots. Independent square classes give degree $2^L$, and all radicands are squares in $\mathbb Q_2$, so 2 splits completely. In disjoint groups of twelve norm-2 prime ideals apply the same exchange. Each group multiplies the density by exactly $1+1/2125764$, since its added mass is $1/16384$ and its original mass is $(3/2)^{12}$. Independent groups multiply the gains, giving the relative factor $(1+1/2125764)^{\lfloor 2^L/12\rfloor}$. This compares each field with its own greedy density.

\subsection{The sharp dimension as a function of field degree}
Order the multiset containing $n$ copies of $1/p$ for each rational prime $p$ in decreasing order and let $S_n(d)$ be the sum of its first $d$ terms.
Among all number fields of degree $n$, the least number of prime-ideal coordinates in which a two-deletion edit strictly improves the greedy cube is
\[
 d_n=\min\{d:S_n(d)>11/2\}.
\]
Every field of degree $n$ requires at least this many coordinates for such an edit. At least one field of degree $n$ attains this dimension. The minimum is finite by divergence of the prime reciprocal sum. The invariant measures exactly the class of edits classified above.

Indeed, for marker weight $w$ the edge formula requires $w^2(1+\sum_i u_i)>1+w$ for a strict gain, since $u_i\le 1/2$ rules out gain from a single spoke. Each satellite contributes $u_i\le w_i$, with equality at satellite bit one. Thus the sum of all the prime-ideal weights must exceed $w+(1+w)/w^2-1$, a decreasing function on $(0,1/2]$ with minimum $11/2$. There are at most $n$ primes above each rational $p$, each of norm at least $p$. Consequently $S_n(d)$ dominates the sum of any $d$ weights in the field.

For the converse, choose a degree-$n$ field in which every rational prime occurring in the first $d_n$ multiset terms splits completely. Such a field exists by coefficientwise Chinese remaindering. Set $g(X)=\prod_{j=1}^n(X-j)$ and, at every specified prime $p$, demand that a monic degree-$n$ polynomial $f$ be congruent to $g$ modulo $p^{2M+1}$ with $M=\max_j v_p(g'(j))$. Hensel's lemma at each $j$ gives a root whose difference from $j$ has valuation $>M$, and the resulting roots are distinct since $v_p(j-k)\le M$ for $j\ne k$. Impose Eisenstein's criterion at one further prime to guarantee irreducibility. Then $\mathbb Q[X]/(f)$ has the required degree and splitting by its tensor products with $\mathbb Q_p$. In these coordinates choose a norm-two marker and take all satellite bits equal to one. The sum condition is sufficient.

\begin{center}
\begin{tabular}{c|rrrrrrrrrr}
$n$&2&3&4&5&6&7&8&9&10&$\ge11$\\ \hline
$d_n$&30\,948&85&27&18&15&14&13&13&12&12
\end{tabular}
\end{center}
For $n\ge12$ the first twelve terms are all $1/2$. For $n=2$ the last prime is $169751$. Integer interval certificates and exact short sums are reproduced by \texttt{check\_field\_thresholds.py}. The case $n=1$ is covered by the formula without requiring enumeration of its prime prefix.

Over the rational integers the threshold nevertheless has a narrow certified scale,
\[
                   3.46\cdot10^{79}<d_1<3.52\cdot10^{79}.
\]
Indeed put $x_-=6.51\cdot10^{81}$ and $x_+=6.57\cdot10^{81}$. The absolute error in replacing the prime reciprocal sum through $x$ by $\log\log x+B$ is at most $1/(2\log^2 x)$ at both endpoints, by Theorem~5 of J.~B.~Rosser and L.~Schoenfeld, \href{https://doi.org/10.1215/ijm/1255631807}{\emph{Approximate formulas for some functions of prime numbers}} (1962). Here $0.261496<B<0.261499$ for the Meissel--Mertens constant. Substitution places the two prime sums on opposite sides of $11/2$. Theorem~1 of the same paper bounds $\pi(x)$ below and above by $(x/\log x)(1+c/\log x)$ with $c=1/2$ and $3/2$, respectively, in this range. Applying these bounds to $\pi(x_-)<d_1\le\pi(x_+)$ gives the interval. The accompanying checker also bounds $B$ directly by truncating $B=\gamma-\sum_p\sum_{k\ge2}1/(kp^k)$ with explicit tail errors.

\section{How many primes must a density rule couple?}
Let $\alpha$ be the supremum of natural densities of subsets of
$\mathbb N_{>0}$ containing no three distinct integers $a,b,c$ with
$ac=b^2$. Partition the primes into blocks of size at most $d$. For each
block choose an arbitrary set of allowed exponent vectors, and admit
an integer when its vector belongs to the allowed set in every block.
Let $\alpha_d$ be the corresponding supremum over all partitions and
rules giving a progression-free set with natural density.
The partition can vary with the competitor. In particular this model
can impose dependence among infinitely many separate groups of primes.

\medskip\noindent\textbf{Theorem (bounded prime blocks).}
For each $d\ge1$, the number $\alpha_d$ is computable and is attained.
An attaining construction uses Rankin's independent exponent rule
outside a finite set of primes. Allowing upper density gives the same
$\alpha_d$ in this class. Furthermore
\[
                  \alpha_d<\alpha,\qquad \lim_{d\to\infty}\alpha_d=\alpha .
\]
Thus every fixed bound on the number of jointly tested primes has a
strict density cost, even after optimizing the partition.
The result does not determine the unrestricted value $\alpha$ or the
first block size that improves on the one-prime rule.

\paragraph{Local probabilities.}
For a finite prime set $B$, the $B$-smooth positive integers form a
semigroup $S_B$. For $F\subset S_B$ write
\[
  c_B=\prod_{p\in B}(1-1/p),\qquad
  \mu_B(F)=c_B\sum_{s\in F}1/s.
\]
This is the natural density of integers with $B$-smooth part in $F$.
Indeed the finite exponent frequencies follow by the Chinese remainder
theorem, and discarding exponents $\ge m$ costs at most
$\sum_{p\in B}p^{-m}$ in upper density. Disjoint finite blocks have
independent probabilities. A nonempty progression-free product forces
each block rule to be progression-free: replace the smooth part of
one admitted integer by the three local terms to see this. Conversely
a product of progression-free rules is progression-free by projection.

Let $G$ consist of the nonnegative integers whose base-three digits
are zero or one, and let $C_B$ require all exponents in $G$. In an
arithmetic progression of three exponents in $G$, the lowest varying
digit would give a nonconstant progression in $\{0,1\}$ modulo three,
which is impossible. Thus $C_B$ is progression-free.
For infinite collections of Rankin primes the product of local
probabilities gives the natural density: the possible violations at
$p>Y$ require $p^2\mid n$ and cost at most $\sum_{p>Y}p^{-2}$.
The infinite product is positive since every factor is at least
$1-p^{-2}$.

\paragraph{Most blocks admit Rankin's rule as an optimizer.}
If $1/c_B\le2$, then $\mu_B(F)\le\mu_B(C_B)$ for every
progression-free $F$. This is the light-block observation following
McNew's Theorem~6.1. Here is a direct proof. For $s\in F\setminus C_B$
replace every ternary digit $2$ in its exponents by $0$, respectively
$1$, to get integers $a,b\in C_B$ with $a<b<s$ and $as=b^2$.
At least one of these proper divisors is missing from $F$. Charge
$1/s$ to one such divisor $t$. Each charged $s$ is a different
$t u$ with $u\in S_B\setminus\{1\}$. The charge at $t$ is at most
$(1/c_B-1)/t\le1/t$. Summing gives the claim.

Call the remaining blocks heavy. If every prime of $B$ exceeds $4d$,
then $\log(1/c_B)\le\sum_{p\in B}1/(p-1)\le1/4<\log 2$.
Hence at most $\pi(4d)$ blocks are heavy, and their union $H$
has at most $h=d\pi(4d)$ primes. Replace every light block by
Rankin's rule. The new set has natural density equal to the finite
heavy probability times the Rankin product off $H$. Its density is
at least the original upper density, since finite collections of
old block restrictions bound the latter from above, and the product
of the corresponding new probabilities converges to the new density.
A heavy rule may likewise be optimized separately. Its maximum exists:
indicators of a maximizing sequence on the countable semigroup admit
a pointwise convergent subsequence by diagonal extraction, avoidance
passes to the limit, and the weights $1/s$ are summable.

\paragraph{Existence of an optimal partition.}
Take a sequence approaching $\alpha_d$, already in the form just
described. Pass to a subsequence with constant number of exceptional
primes and constant partition of their ordered positions. Each of the
at most $h$ prime labels, after another subsequence, is either constant
or tends to infinity. In every block temporarily impose valuation zero
at each escaping prime $p$. The density loss is at most the sum of
$1/p$ over these primes, hence tends to zero. The surviving local
zero slices are progression-free on the fixed prime labels. One
may now use the optimal rule on each fixed block and Rankin's rule
at each discarded prime, since zero is an allowed Rankin exponent.
This single fixed construction attains at least the limiting density,
and belongs to the original class. The upper-density assertion follows
from the replacement bound above.

\paragraph{A uniform improvement.}
Fix the at most $h$ exceptional primes of a competitor after replacement.
A marker prime $\ell$ outside $H$ can be taken at most $L$, the
$(h+1)$st prime. Choose $Y\ge L$, depending only on $d$, so large that
\[
          \sum_{q\le Y,\ q\ {\rm prime}}\frac1q
                > 2L^2+(h+1)/2 .
\]
This is possible by Euler's divergence theorem. For the satellite set
$B$ use all primes through $Y$ outside $H\cup\{\ell\}$. Thus
$\sum_{q\in B}1/q>2L^2$. Apply the exchange of Section~\ref{sec:cap} to the lowest digits at the marker and satellites. Its cap proof works with any marker and leaves the higher digits binary.

Write $w=1/\ell$, $u_q=1/q$. For fixed good higher digits the change
in unnormalized reciprocal mass, relative to their common weight, is
\[
   \left(\prod_{q\in B}u_q\right)
       \left[w^2\left(1+\sum_{q\in B}u_q\right)-(1+w)\right]
      >\tfrac12\prod_{q\in B}u_q .
\]
The inequality uses $w\ge1/L$, $w\le1/2$. The higher digits have
total weight at least one. With $P_Y=\prod_{p\le Y}p$ and
$c_Y=\prod_{p\le Y}(1-1/p)$, the new local probability exceeds the
old one by at least $\delta_d=c_Y/(2P_Y)>0$. Since local probability
is at most one, this increases the full product by a factor at least
$1+\delta_d$. Thus $\alpha\ge(1+\delta_d)\alpha_d>\alpha_d$.
In particular, even allowing each partition to use distant, arbitrarily
large primes gives a uniform gap for fixed $d$.

\paragraph{Approximation by a finite block.}
Let $P$ be the primes through $Y\ge2$, let $c=c_P$ and let $M$ be the
largest reciprocal mass $\sum_{s\in F}1/s$ for a progression-free
$F\subset S_P$. For any progression-free set $A$ of natural density,
write $n=rs$ with $(r,\prod_{p\in P}p)=1$ and $s\in S_P$.
For each fixed $r$ the permitted $s$ form a progression-free set.
Consequently
\[
 \sum_{n\le X,\,n\in A}\frac1n
  \le M\sum_{r\le X,\,(r,\prod_{p\in P}p)=1}\frac1r
  =M(c\log X+O_Y(1)).
\]
The last estimate follows by finite inclusion--exclusion; partial
summation gives $d(A)$ as the limit of the left side divided by
$\log X$. This proves $\alpha\le cM$. A maximizing smooth template
on $P$ and Rankin's rule elsewhere, which has width $\pi(Y)$, gives
\[
  \alpha_{\pi(Y)}\ge cM\prod_{p>Y}(1-p^{-2})
       \ge cM\,\frac{Y}{Y+1}
\]
for integer $Y$, by telescoping over all integers greater than $Y$.
Together with monotonicity this proves convergence. Choosing $Y=p_{d+1}-1$, where $p_j$ is the $j$th prime, gives $0<\alpha-\alpha_d\le 1/p_{d+1}$ since $cM\le1$. Thus $\alpha$ itself is computable.

\paragraph{Effective computation for fixed width.}
Uniformly over competitors after replacement, freezing exponents at
heavy primes $p>Y$ to zero costs at most $h/Y$. Releasing just
these primes to use Rankin's rule can only increase density. Thus
optimizing partitions of primes through $Y$ of width $d$, with the
Rankin tail, approximates $\alpha_d$ from below with error at most
$h/Y$. There are finitely many partitions to search. Truncate each
prime exponent below $m$; in total this costs at most
$\sum_{p\le Y}p^{-m}$. On finite boxes, enumerate all subsets,
reject any nonconstant arithmetic progression of exponent vectors,
and maximize the exact rational probability. Finally the Rankin
tail product through any $Z>Y$ has an absolute truncation error at most
$\sum_{n>Z}n^{-2}\le1/Z$; its single-prime factors can be truncated
geometrically. These estimates give rational intervals of arbitrary
prescribed width. They make no assertion of practical running time.

\paragraph{Relation to earlier work.}
McNew's Theorems 5.1--5.2 concern unrestricted \emph{upper} density,
and his Theorem~6.1 treats a one-prime logarithmic optimum.
See \href{https://www.nathanmcnew.com/GPFsets.pdf}{McNew's paper,
Sections 5--7}. The incoming manuscript
\texttt{mcnew\_counterexample\_and\_grh\_transfers.zip} disproves
Conjecture~7.2 by a marker and a shared deletion and shows that any
fixed finite change to Rankin's rule admits improvement. The bounded-block
theorem allows an arbitrary rule in every one of infinitely many
blocks, proves a gap uniform over all their partitions, and gives
attainment and effective approximation for each width. The lowest-digit
construction above supplies its local gain.

\subsection{The exact optimum through block size four}
Write $R$ for Rankin's set. In the notation of the preceding theorem,
\[
       \alpha_1=\alpha_2=\alpha_3=\alpha_4
       =d(R)=\prod_p\left[(1-1/p)\prod_{j\ge0}(1+p^{-3^j})\right],
\]
where the outer product runs over primes.
In particular, any progression-free product rule of density strictly
greater than $d(R)$ must couple at least five primes in some block.
An incoming manuscript had established local optimality at the single
four-tuple $(2,3,5,7)$. To pass from one tuple to all tuples we certify
a family of rectangular prefixes, then change the four prime weights by
summation by parts. This quantifier change matters because a product
partition can pair even the smallest prime with an arbitrarily large one.

It suffices to prove optimality on each four-prime block: one can pad a
smaller block with extra primes bearing independent Rankin rules and
divide by their positive local probabilities. Order the primes as
$q_1<q_2<q_3<q_4$. If $q_1\ge5$, then
$\prod q_i/(q_i-1)\le(5/4)(7/6)(11/10)(13/12)<2$, and the light-block
bound applies. Otherwise select the following base tuple $b$.
The first $f$ entries are fixed at $q_i=b_i$; all others obey $q_i\ge b_i$.
\begin{center}\small
\begin{tabular}{lrrrr}
Condition & Base tuple & $f$ & Prefixes & Search nodes\\ \hline
$q_1=2,\ q_2=3$ & $(2,3,5,7)$ & 2 & 9 & 5,537,385\\
$q_1=2,\ q_2\ge5$ & $(2,5,7,11)$ & 1 & 81 & 221,071\\
$q_1=3$ & $(3,5,7,11)$ & 1 & 81 & 29,307
\end{tabular}
\end{center}

Use the box $D=\prod_i\{0,\ldots,m_i-1\}$ with
$m=(9,9,3,3)$, corresponding to low ternary digit depths
$k=(2,2,1,1)$. Write $G^4$ for the good vectors,
$W_b(J)=\sum_{a\in J}\prod_i b_i^{-a_i}$ and $\lambda=21/20$.
For every rectangular prefix $D'$ of $D$ with the first $f$ axes
untruncated, the program \texttt{check\_four\_prime\_blocks.py}
certifies that any progression-free $I\subset D'$ satisfies
\[
       \lambda W_b(I\setminus G^4)\le W_b((G^4\cap D')\setminus I).
       \tag{*}
\]
The test orders the box by $n=\prod b_i^{a_i}$ and branches on including
or omitting the least undecided value. Inclusion of the midpoint $v$,
with an already chosen $u<v$, forbids $v^2/u$ whenever this is in the
box. The upper bound on the objective is the summed weight of all chosen
and undecided values, with the factor $\lambda$ on bad values. A branch
is complete as soon as this bound is at most $W_b(G^4\cap D')$.
The objective is scaled to integers. Every independent subset follows
one branch, proving that a successful traversal certifies the inequality.

For an arbitrary ratio $t_i=b_i/q_i\le1$, the function
$t_i^{a_i}$ on $0\le a_i<m_i$ is a nonnegative linear combination
of prefix indicators. The coefficient of $\mathbf1_{a_i\le j}$ is
$t_i^j-t_i^{j+1}$ for $j<m_i-1$ and $t_i^{m_i-1}$ for the last
prefix. A fixed coordinate uses just the last prefix.
Multiply the four expansions, apply (*) to $I\cap D'$ on each prefix
and sum. This gives (*) on the full box with every $b_i$ replaced by
$q_i$.

For completeness, the low-box margin propagates to unbounded exponents
as follows. Let $F$ be a progression-free subset of the four-prime
semigroup and put $M=W_q(C\setminus F)$, where $C$ is its Rankin set.
Split exponents into low digits and a high part divisible in each
coordinate by $m_i$. In each fixed good high part, (*)
bounds the mass of points in $F$ with a bad low digit by the missing
good mass divided by $\lambda$. Summing gives at most $M/\lambda$.
For a point in $F$ with a bad high digit, replace every digit two by
zero, respectively one. As in the light-block proof, one of the two
resulting Rankin points must be missing. Relative to a missing point
the charged reciprocal weight is a product of factors $q_i^{-e3^j}$
for a subset of digits, with $e=1$ for the one replacement and $e=2$
for the zero replacement. At least one changed digit is high.
The total charge is at most $M T(q)$, where
\[
 T(q)=\sum_{e=1,2}\left[\prod_i\prod_{0\le j<k_i}
       (1+q_i^{-e3^j})\right]
       \left[\prod_i(1-q_i^{-e m_i})^{-1}-1\right].
\]
The second bracket bounds all nonempty high-digit changes by geometric
series; ternary subset sums have distinct exponents. The bound decreases
when a prime increases. For each base tuple the program checks
$1-1/\lambda-T(b)>1/200$ in rational arithmetic. Therefore
$W_q(C)-W_q(F)\ge M(1-1/\lambda-T(q))\ge0$.
This proves the claimed local optimum and, by the finite-block
upper bounds proved above, the global result. The base-box
branching and high-digit charge adapt the incoming
\incoming{mcnew_density_and_grh_transfers.zip}{four-prime density manuscript}. Its three
fixed-prime certificates alone do not imply the prefix inequalities.

\section{Sharp precision for intersecting Lubin--Tate branches}
Let $K/\mathbb Q_p$ be finite, $\mathcal O$ its ring of integers, $\pi$ a uniformizer and $q$ the size of its residue field. Normalize $v(\pi)=1$. For the Lubin--Tate formal $\mathcal O$-module attached to $\pi$, write $[a](X)$ for multiplication by $a$. Choose a finite set $B\subset\mathcal O$ of $d\ge2$ distinct directions. Fix $b\in B$ and put
\[
 \Delta_b=\prod_{c\in B\setminus\{b\}}([b](X)-[c](X)),\quad
 h_c=v(b-c),\quad \omega=\sum_{c\ne b}q^{h_c}.
\]
Over $A_r=\mathcal O/\pi^r$, $r\ge1$, define $T_r$ to be the least $T$ for which $X^T/\Delta_b\in A_r[[X]]$. Inversion takes place in $A_r((X))$. A distinguished factor $X^s+\pi V$ is invertible there by a geometric expansion of length $r$.

The earlier \emph{Cycle signings and Lubin--Tate} report \cite{priorlt}, \texttt{lubin\_tate.tex}, uses this pole to measure the precision of an observation. At each $c$ one observes $H(X,[c](X))\bmod X^{n_c}$ for unknown $H\in A_r[[X,Y]]$. The componentwise least precision at $b$ that permits recovery of every coefficient of total degree $<d$ is $1+T_r$, with the corresponding minimum at each other branch. Indeed, divide $H$ by the monic polynomial $\prod_c(Y-[c](X))$ using formal Weierstrass division. Its terms have total degree $\ge d$, so the remainder has the desired jet. Lagrange interpolation of the remainder multiplies each observed value by $P_b/\Delta_b$, where $P_b=\prod_{c\ne b}(Y-[c](X))$. An error divisible by $X^{1+T_r}$ therefore only affects terms in $(X,Y)^d$. Conversely $(X^{T_r}/\Delta_b)P_b$ vanishes on the other branches, has value $X^{T_r}$ at $b$, and has nonzero constant-in-$X$ coefficient at $Y^{d-1}$. The earlier report gives the bound $T_r\le\omega+(r-1)(q^{\max h_c}-1)$ and asks in its third further research question when simultaneous branches attain it.

\paragraph{Theorem.}
If all $h_c=0$, then $T_r=\omega$ for every $r$. Otherwise write
\[
 t=\max h_c,\qquad D=q^t-q^{t-1},\qquad m=\#\{c:h_c=t\},\qquad U_r=\omega+(r-1)D .
\]
Then
\[
 U_r-(m-1)D\ \le\ T_r\ \le\ U_r,\qquad
 T_r=U_r\ \Longleftrightarrow\
 \binom{m+r-2}{r-1}\not\equiv0\pmod p .
\]
Equivalently, equality holds precisely when adding $r-1$ to $m-1$ produces no carry in base $p$. In particular the pole grows with exact slope $D$. When the deepest contact is unique the displayed upper bound is exact at every precision.

\paragraph{Proof.}
The difference $[b]-[c]$ is the product of $[b-c]$ and a unit, since $F(U,V)-U$ is $V$ times a unit for the formal group law $F$. Also $[b-c]=[a]\circ[\pi]^{\circ h_c}$ with $a$ a unit, and $[a](Z)/Z$ is a unit. Multiplying by a unit does not change the pole. An invertible coordinate substitution $X\mapsto X u(X)$ also preserves it since it preserves every fractional lattice $X^jA_r[[X]]$. By the Lubin--Tate isomorphism theorem for one uniformizer \cite{lt}, we may take $[\pi](X)=\pi X+X^q$. Set $f_0=X$ and $f_h=\pi f_{h-1}+f_{h-1}^q$. We have reduced the full denominator to the monic polynomial $\prod_c f_{h_c}$. This proves that the contact multiset completely determines $T_r$ for the fixed field and uniformizer.

For $h\ge1$ write $s=q^h$, $D_h=s-q^{h-1}$. Then
\[
 f_h(X)\equiv X^s+\pi X^{s-D_h}\pmod{\pi^2}.
\]
Indeed $f_{h-1}\equiv X^{q^{h-1}}\bmod\pi$ and replacing it by this monomial in the $q$th power is valid modulo $\pi^2$ because $\pi\mid q$. The polynomial $X^{-s}f_h$ therefore has the form $1+\pi X^{-D_h}+\pi^2 V_h(X^{-1})$, with $\deg V_h\le s-1<2D_h$, since $f_h(0)=0$. For $h=0$ the normalized factor equals one.

Put $z=X^{-1}$. Give a term of $\pi$-adic order $j$ a degree allowance $jD$. In the finite nilpotent expansion of $X^{\omega}/\prod_c f_{h_c}$ modulo $\pi^r$, every positive-order term obeys this allowance. A factor with $h<t$ has strict inequality, as does every contribution from $\pi^2V_h$. To reach degree $(r-1)D$ one must spend exactly $r-1$ powers of $\pi$ on the terms $\pi z^D$ from the $m$ deepest factors. Its coefficient is thus congruent modulo $\pi^r$ to
\[
 (-1)^{r-1}\binom{m+r-2}{r-1}\pi^{r-1},
\]
the corresponding term of $(1+\pi z^D)^{-m}$. All greater degrees vanish. This proves the upper bound and its equality test. If $t=0$, each normalized factor is one.

For the lower bound put $a_k=\binom{m+k-1}{k}\bmod p$. Its generating series is $(1-z)^{-m}$. The recurrence from multiplication by $(1-z)^m$ and the nonzero constant and leading coefficients forbid $m$ consecutive zeros: such a block would force all later coefficients to vanish, while $(1-z)^{-m}$ is not a polynomial. Since $a_0=1$, some $j$ in $[\max(0,r-m),r-1]$ has $a_j\ne0$. Reduction modulo $\pi^{j+1}$ gives $T_r\ge T_{j+1}=\omega+jD$. The carry formulation follows from the base-$p$ binomial criterion.\hfill$\square$

\paragraph{Answer to the earlier question.}
At $r=1$ the original bound is always attained. For $r>1$ with $t=1$, its slope $q-1$ equals $D$ and the binomial criterion classifies all cases. If $t\ge2$, then $D<q^t-1$, so no branch attains the original bound. When all contacts have exponent zero the original bound is exact. For example in $\mathbb Q_2$ take $B=\{0,4,12\}$ and $b=0$. Here $\omega=8$, $D=2$ and $m=2$. The pole is $8$ at $r=2$, where $f_2^2\equiv X^8\bmod4$, and $12$ at $r=3$, where the top coefficient survives. The corresponding observation precisions are $9$ and $13$. These are rigorous finite-level consequences of the interpolation problem and require no zero-location hypothesis.

\paragraph{Corollary (optimal sampling at every modulus).}
Suppose $d\ge2$ directions must be chosen inside one coset of $\pi^a\mathcal O$, $a\ge0$. For $1\le j<d$ put
\[
 H_j=a+\max\{h:q^h\mid j\}.
\]
The least possible common observation precision is $1+T_r(H_1,\ldots,H_{d-1})$, where the notation means the pole of the reciprocal of $\prod_{j=1}^{d-1} f_{H_j}$. A single set of $d$ directions attains this simultaneously for every $r$. Choose representatives of the residue field numbered $0$ through $q-1$ with $0$ representing zero. Map the base-$q$ digits of $0,\ldots,d-1$ to these representatives in a $\pi$-adic expansion, then scale by $\pi^a$ and translate into the required coset.

To see optimality, first observe that $f_u$ divides $f_v$ in $\mathcal O[X]$ for $u\le v$, by iteration and $f_1(0)=0$. Replacing any contact exponent with a larger one therefore cannot decrease the reciprocal pole: multiply a power series clearing the larger denominator by the integral quotient to clear the smaller denominator. Scale the coset back to $\mathcal O$ for the counting argument. Starting with any $d$ directions, continue down the most populous child residue class. Its occupancy at depth $h$ is at least $\lceil d/q^h\rceil$, by induction. The path selects some $b$ since the finite set is nonempty at every depth. The number of other directions with $v(b-c)\ge h$ is at least $\lceil d/q^h\rceil-1=\lfloor(d-1)/q^h\rfloor$. These are the tail counts of the list $H_j-a$. Sorted in decreasing order, this branch's contacts dominate that list coordinatewise. Conversely in the indicated digit construction each residue class at depth $h$ has at most $\lceil d/q^h\rceil$ members, so every branch's sorted contacts are bounded above by the list. The zero-indexed direction reaches equality at every depth. Restoring the shift $a$ and applying monotonicity proves the claim.

For example $d=q^k+1$ with $k\ge1$ has a unique deepest contact in the extremal list. Counting the multiples of $q^h$ gives $\omega=q^a(q^k+k(q-1)q^{k-1})$, so the optimal common precision is exactly
\[
  1+q^a\big(q^k+k(q-1)q^{k-1}\big)+(r-1)(q-1)q^{a+k-1}.
\]
For $q=2$, five directions constrained to one coset of $\pi\mathcal O$ need $17+4(r-1)$ coefficients per branch.

\subsection{The exact defect at every level}
Keep the notation above and assume $t\ge1$. Put $d_0=q-1$ and $a=D/d_0=q^{t-1}$. All exponents occurring in $f_h(X)$ are $1\bmod d_0$, by induction in the recurrence for $f_h$. Consequently the reciprocal polynomial satisfies
\[
   X^{-\omega}\prod_{c\ne b} f_{h_c}(X)=F(X^{-d_0}),\qquad F\in\mathcal O[z],\quad F(0)=1 .
\]
In particular the loss $U_r-T_r$ is always divisible by $d_0$. Write $v_p$ for the ordinary $p$-adic valuation on integers, and set $e=v(p)$.

\paragraph{Theorem.}
There exists a bounded eventually periodic sequence $\eta_r$ of integers such that
\[
 T_r=\omega+d_0\{a(r-1)-\eta_r\},\qquad 0\le\eta_r\le (m-1)a.
\]
The following construction computes its exact values, including the initial ones, from a single finite ring independent of $r$.
Adjoin $\rho$ with $\rho^a=\pi$, put $N=ma$ and use
$S=\mathcal O[\rho]/(\rho^N)$. Here $v_\rho(\rho)=1$, with the convention that the valuation of a zero class in $S$ is $+\infty$. Set
\[
 G(z)=F(z/\rho),\qquad
 \frac1{G(z)}=\sum_{j\ge0}g_jz^j\quad\text{in }S[[z]],\qquad g_j=0\quad(j<0).
\]
Then $\eta_r$ is the least integer $i\in[0,(m-1)a]$ satisfying
\[
             v_\rho\bigl(g_{a(r-1)-i}\bigr)<a+i .
\]
\paragraph{Proof.}
The degree allowance established above gives $a v([z^j]F)\ge j$, so $G$ is integral. On reducing $G$ modulo $\rho$, only the constant terms and the first-order terms from the $m$ deepest factors survive. Thus $G\equiv(1+z^a)^m\bmod\rho$. If $H_j=[z^j]F^{-1}$, then $g_j$ is the reduction of $H_j\rho^{-j}$. At $j=a(r-1)-i\ge0$ the condition $v(H_j)<r$ is exactly $v_\rho(H_j\rho^{-j})<a+i$, since $v(H_j)$ is integral. All these cutoffs are at most $N$. The preceding theorem ensures there is a surviving $H_j$ in this range. The largest surviving $j$ gives the pole $\omega+d_0 j$, proving the formula.

Here is an explicit, possibly nonminimal, period. Put $L=mN$, $E=G-(1+z^a)^m$ in $S[z]$, and
\[
 C(z)=\sum_{k=0}^{N-1}(-E)^k(1+z^a)^{m(N-1-k)} .
\]
Since $E^N=0$, the identity $G^{-1}=C(z)/(1+z^a)^L$ holds in $S[[z]]$. For each class modulo $a$, and for $j\ge\deg C$, the coefficient $g_j$ is a finite linear combination over $S$ of shifted sequences $(-1)^u\binom{u+L-1}{L-1}$, indexed by $u=(j-k)/a$ for the relevant monomials $z^k$ in $C$. Let
\[
 b_0=\lceil m/e\rceil,\qquad
 A=2p^{\,b_0+\lfloor\log_p(\max\{1,L-1\})\rfloor}.
\]
Then $p^{b_0}=0$ in $S$ because $v_\rho(p)=ae$. For $1\le l\le L-1$, the identity $l\binom A l=A\binom{A-1}{l-1}$ gives $p^{b_0}\mid\binom A l$. Vandermonde's identity shows that the displayed signed binomial sequence has period dividing $A$ in $S$. This includes $L=1$, when evenness of $A$ suffices. Hence $g_{j+aA}=g_j$ for $j\ge\deg C$, and $\eta_{r+A}=\eta_r$ whenever $a(r-1)-(m-1)a\ge\deg C$.\hfill$\square$

\paragraph{A closed formula at depth one.}
If $t=1$, then $F(z)=(1+\pi z)^m$. Here the full answer simplifies to
\[
 T_r=\omega+(q-1)\max\left\{k\in\{0,\ldots,r-1\}:         k+e\,v_p\binom{k+m-1}{m-1}<r\right\}.
\]
Indeed the $k$th coefficient of $F^{-1}$ is $(-\pi)^k\binom{k+m-1}{m-1}$. For example, with $q=p=2$, $m=2$ and $e=1$, the maximum is $r-1$ for odd $r$ and $r-2$ for even $r$. One realizes this example at $b=0$ with directions $0,2,6$ over $\mathbb Z_2$.

\paragraph{Computation over the original local ring.}
For implementation the auxiliary root $\rho$ can be removed. Define $k_j=\lceil j/a\rceil$, $c_i=[z^i]F/\pi^{k_i}$ and $B_j=H_j/\pi^{k_j}$. Dividing the recurrence $[z^j](FF^{-1})=0$ by $\pi^{k_j}$ gives $B_0=1$ and, for $j>0$,
\[
 B_j=-\sum_{i=1}^{\min(j,\deg F)}c_i B_{j-i}\pi^{k_i+k_{j-i}-k_j}\pmod{\pi^m}.
\]
The exponents in this expression are zero or one. For $j=a(r-1)-i\ge0$, the survival test becomes $B_j\not\equiv0\bmod\pi^{\lfloor i/a\rfloor+1}$. This recurrence computes every level in the fixed ring $\mathcal O/\pi^m$.

At precision $r=2$ there is also a formula for arbitrary depths. Let $m_h=\#\{c\ne b:h_c=h\}$. The congruence for $f_h$ modulo $\pi^2$ gives
\[
 T_2=\omega+\max\bigl(\{0\}\cup
   \{q^h-q^{h-1}:h\ge1,\ p\nmid m_h\}\bigr).
\]
For $p=2$ and $q=4$, compare contact multisets $(0,2,2)$ and $(1,2,2)$. Both have $D=12$ and $m=2$, and neither reaches $U_2$. In the first case the loss is $12$. In the second the depth-one factor leaves a surviving pole increase of $3$, and the loss is $9$. Thus shallower contacts can determine the exact answer even at the first cancellation.

\paragraph{A classical special case.}
Over $K=\mathbb Q_p$ take $\pi=p$ and a single factor of depth $t\ge1$. In the multiplicative formal group its division polynomial is $(1+X)^{p^t}-1$. The condition $X^n/\Delta_b\in A_r[[X]]$ says exactly that $X^n=0$ in the quotient by this polynomial. Weierstrass division identifies the quotient with $(\mathbb Z/p^r)[T]/(T^{p^t}-1)$ via $X=T-1$. Our value $p^t+(r-1)(p-1)p^{t-1}$ agrees with the classical nilpotency index of the augmentation ideal of this cyclic group ring. See P. L. Clark and U. Schauz, \emph{Functional degrees and arithmetic applications, I}, Theorem~4.7 and its second proof, which recalls results of Weisman and Wilson (\url{https://arxiv.org/abs/2201.02763}). The multibranch problem involves products of division polynomials in the same variable, including repeated factors.

\section{A simultaneous congruence for weighted Fleck sums}
The multiple-branch calculation has an arithmetic interpretation involving only ordinary integers. Fix a prime $p$, a modulus $M=p^t$ with $t\ge1$, and a positive integer $m$. For $0\le u<M$ and $0\le j<m$ define
\[
 S_{j,u}(n)=\sum_{\substack{0\le k\le n\\ k\equiv u\bmod M}}
 (-1)^{n-k}\binom nk\binom{(k-u)/M}{j}.
\]
The lower argument $(k-u)/M$ is nonnegative at every summand. Sums of this type are the cyclotomic $\psi$-coefficients before normalization, up to the displayed sign. Sun and Wan prove Lucas-type congruences for individual normalized coefficients in \emph{Lucas-type congruences for cyclotomic $\psi$-coefficients} (\url{https://arxiv.org/abs/math/0512012}). Their integrality statement gives $v_p(S_{j,u}(n))\ge\lfloor(n-M/p-jM)/(M-M/p)\rfloor$. At $j=m-1$ it supplies the sufficient common threshold $U_r$ below. Here we determine the least simultaneous threshold and identify the obstruction one step earlier. The derivation is independent of any claim to priority about these congruences.

\paragraph{Corollary.}
For $r\ge1$ take $T_r$ from the Lubin--Tate calculation with $K=\mathbb Q_p$, $\pi=p$ and the contact list consisting of $m$ copies of $t$. Then
\[
  p^r\mid S_{j,u}(n)\quad\text{for every }j,u
  \quad\Longleftrightarrow\quad n\ge T_r.
\]
Here $\omega=mM$, $D=M-M/p$, and $U_r=mM+(r-1)D$. Thus
$U_r-(m-1)D\le T_r\le U_r$, with equality at the upper endpoint precisely when $r-1$ and $m-1$ add without a base-$p$ carry. For every $r$, the values $S_{m-1,u}(T_r-1)$ are all equal modulo $p^r$ to a single nonzero class, and $S_{j,u}(T_r-1)\equiv0$ for $j<m-1$. When the upper endpoint is attained, the common class is
\[
            (-p)^{r-1}\binom{m+r-2}{r-1}\pmod {p^r}.
\]
\paragraph{Proof.}
Set $R=\mathbb Z/p^r$, $W=Y^M-1$, $P=W^m$. In the multiplicative Lubin--Tate coordinate $X=Y-1$ the denominator is associated to $P$. Hence $T_r$ is the nilpotency index of $Y-1$ in $R[Y]/P$, by monic division as in the cyclic special case above. The polynomials $Y^uW^j$ with $0\le u<M$, $0\le j<m$, form a basis over $R$. Their leading monomials run through degrees $0$ to $mM-1$ with coefficient one. Writing $Y^{u+Ml}=Y^u(1+W)^l$ in the binomial expansion of $(Y-1)^n$ identifies its coordinates in this basis with $S_{j,u}(n)$. This proves the threshold.

Let $Q(Y)$ be the polynomial remainder of $(Y-1)^{T_r-1}$ on division by $P$. Then $Q\ne0$ and $(Y-1)Q= cP$ for some scalar $c\in R$, by the degree bound and the vanishing at $T_r$. Multiplication by the monic $Y-1$ is injective in $R[Y]$, whence
\[
 Q=cW^{m-1}(1+Y+\cdots+Y^{M-1}),\qquad c\ne0 .
\]
This gives the claimed full vector of residues at the last exceptional index.

For the explicit value of $c$ at saturation, put $P_X=((1+X)^M-1)^m$ and $F_*(z)=z^{mM}P_X(1/z)$. In the variable $X$ the leading coefficient in degree $mM-1$ of the remainder of $X^{T_r-1}$ is still $c$. Monic division at infinity identifies it with $[z^{T_r-mM}]F_*^{-1}$, by comparing the coefficients of $X^{-1}$ in $X^{T_r-1}/P_X$. In the reverse polynomial of a single factor the nonconstant coefficients are $\binom Mj$ in deficit $M-j$, $1\le j<M$. They have valuations $t-v_p(j)$, since $\binom{M-1}{j-1}$ is prime to $p$. The unique coefficient with deficit equal to its valuation times $D$ is at $j=M/p$. Its coefficient divided by $p$ is $1\bmod p$, as follows by reducing $\binom{M-1}{M/p-1}$ modulo $p$. The top-degree computation for $(1+pz^D)^{-m}$ now gives the stated residue.\hfill$\square$

For example, take $p=2$, $t=2$, $m=2$, $r=2$. The general upper endpoint is $10$ and the exact threshold is $8$. Indeed $(Y^4-1)^2\equiv(Y-1)^8\bmod4$ by the binomial theorem. At $n=7$ direct expansion gives $S_{1,u}(7)\equiv1\bmod4$ for every $u$. The direct integer check in the archive computes these residues, as well as both sides of the threshold for further parameters.

\section{What remains open}
\paragraph{Density.}
The counterexample settles the proposed equality $\alpha=d(R)$. The proofs give computable approximations to $\alpha$. Its exact numerical value remains undetermined. The smallest width $d$ giving $\alpha_d>d(R)$ lies between $5$ and $225$, inclusive: the four-prime theorem supplies the lower bound, and the 224 odd primes plus the marker in \cite{incomingmc} give the upper bound. The two-deletion theorem identifies the first profitable change in number of deleted residue patterns, a different measure of complexity. Can three or more deletions give a comparably explicit classification, or an appreciable density gain using few ordinary primes? The polynomial construction beats the greedy set for rational-function ratios. The best upper density for \emph{polynomial} ratios in Asada et al.'s Conjecture 6.3 remains a separate target.
\paragraph{Interpolation.}
The saturation criterion uses only the deepest contact and its multiplicity. At nonsaturating levels the other contacts can affect the precise loss. A minimal period in terms of the contacts, or a closed form for the full loss over ramified fields, would improve the finite recurrence. Berger and Sprang's Lemma 4.2 gives a finite residue-field observation in their character-variety argument \cite{bs}; the result here concerns the exact cost of measurements over thicker coefficient rings. Colovi\'c has separately announced the global bounded-functions equality for finite extensions \cite{colovic}.
\paragraph{Formalization.}
The cap argument first needs only a ternary vector space, the third-point map $(x,y)\mapsto-x-y$, and a finite measure. A Lean implementation can prove line exclusion for the explicit replacement before introducing prime valuations or density limits. For the interpolation result, first formalize the reversed iterates $f_h$ over a discrete valuation ring and their weighted degree bound. Then prove the leading coefficient identity over the quotient by $\pi^r$ and the monic-division description of the pole. This isolates two finite algebraic cores that admit comparison with the included exact checkers.
\begin{thebibliography}{99}
\bibitem{mcnew} N. McNew, \emph{On sets of integers which contain no three terms in geometric progression}, Math. Comp. 84 (2015), 2893--2910. \href{https://www.nathanmcnew.com/GPFsets.pdf}{Author's text}, especially Theorem 6.1 and Conjecture 7.2.
\bibitem{incomingmc} ProveIt, \incoming{mcnew_counterexample_and_grh_transfers.zip}{counterexample and GRH transfers}. The files \texttt{rankin.tex} and \texttt{finite\_rankin.tex} prove a strict natural-density improvement by exact-valuation exchanges.
\bibitem{asada} M. Asada, E. Fourakis, S. Manski, N. McNew, S. J. Miller, G. Moreland, \emph{Subsets of $\mathbb F_q[x]$ free of 3-term geometric progressions}, Finite Fields Appl. 44 (2017), 135--147. \href{https://www.nathanmcnew.com/fqxramsey.pdf}{Author's text}. See also the \incoming{grh-progression-densities-and-cubic-prime-race.zip}{incoming polynomial-density report} and \incoming{grh_transfers_density_gaps_and_zero_tests.zip}{density gaps and zero tests}.
\bibitem{best} A. Best et al., \emph{Geometric-progression-free sets over quadratic number fields}. \href{https://arxiv.org/abs/1412.0999}{arXiv:1412.0999}.
\bibitem{idealcount} M. R. Murty and J. Van Order, \emph{Counting integral ideals in a number field}, Expo. Math. 25 (2007), 53--66. \href{https://doi.org/10.1016/j.exmath.2006.08.001}{DOI}.
\bibitem{priorideal} ProveIt, \incoming{grh_ideal_densities_and_transfer_conditions_2026-11-26.zip}{ideal densities and transfer conditions}, sections on integral-ideal density and the $S_3$ transfer problem.
\bibitem{priorlt} ProveIt, \incoming{grh_transfers_cycle_signings_and_lubin_tate.zip}{cycle signings and Lubin--Tate interpolation}, \texttt{lubin\_tate.tex} and Question 3 in \texttt{report.tex}. The present pole calculation answers that question.
\bibitem{lt} J. Lubin and J. Tate, \emph{Formal complex multiplication in local fields}, Ann. of Math. 81 (1965), 380--387. \href{https://doi.org/10.2307/1970622}{DOI}.
\bibitem{bs} L. Berger and J. Sprang, \emph{Integer-valued polynomials and $p$-adic Fourier theory}, \href{https://perso.ens-lyon.fr/laurent.berger/articles/POLINT.pdf}{author's text}.
\bibitem{colovic} U. A. Colovi\'c, \emph{On the ring of bounded functions on the character variety}, 2026 preprint. \href{https://arxiv.org/abs/2611.25162}{arXiv:2611.25162}.
\end{thebibliography}
\end{document}
